\subsection{Statistical Zero-Knowledge of Horizon Interactions: Comb--Diamond Norm Criterion, Decoupling Mechanism, and \texorpdfstring{$\zeta$}{zeta}-Spectral Closed Loop}\label{subsec:discussion-horizon-hszk-comb-sff-zeta}
The previous subsection used ``truncated horizon multipole moments/Toeplitz--PSD certificates'' as observable algebraic constraints, providing an information-theoretic reformulation of no-hair and the distinguishability area law derived from the endpoint quadratic compression (\S\ref{subsec:discussion-horizon-zk-toeplitz-area-law}).
This subsection elevates that caliber to the quantum interactive level: characterizing the horizon interface as a quantum comb (the pluggable skeleton of quantum strategies/quantum networks), and providing a verifiable criterion for ``horizon statistical zero-knowledge'' (HSZK) via the diamond norm of interactive maps; then quantitatively aligning HSZK error with externally extractable information (Holevo leakage), and converting the chaotic scrambling mechanism into zero-knowledge simulability via ``approximate 2-design \(\Rightarrow\) decoupling''; finally providing an operational spectral diagnostic path: controlling the approximate design threshold by the frame potential and spectral form factor (SFF), and in the \(\zeta\)-spectral model identifying the zero-point structure as the unique spectral input determining the time behavior of zero-knowledge error.
\par
The role of this subsection is solely interface docking and falsifiable alignment: all assertions except those explicitly labeled as conjectures are standard quantum information theory inequalities or their direct corollaries, and do not serve as premises of the main proof chain.
\par

\paragraph{Macroscopic charge fiber and horizon comb}
Let \(\mathcal{H}_A\) be the microscopic degrees of freedom of the system (which may be taken to be finite-dimensional), and let \(\Omega\) be its state set (pure or mixed).
Let \(Q:\Omega\to\mathcal{Q}\) be the macroscopic charge map (e.g., discretized sector labels for mass/angular momentum/charge), and denote the charge fiber as \(Q^{-1}(q)\).
For any interaction duration \(t\ge 0\), the horizon interface is given in quantum comb form:
\[
\mathsf{H}_{\le t}:\Omega\to \mathrm{Comb}\bigl(\mathsf{In}_{\le t}\to \mathsf{Out}_{\le t}\bigr),
\]
meaning: an external verifier can adaptively inject probe states during \([0,t]\) and measure the output, while \(\mathsf{H}_{\le t}(\omega)\) gives the equivalence class of all interactive CPTP strategies permitted under microstate \(\omega\) (see \cite{ChiribellaDArianoPerinotti2008QuantumCircuitArchitecture,ChiribellaDArianoPerinotti2009FrameworkQuantumNetworks}).
Let the external strategy (verifier) be any permitted comb \(\mathsf{V}_{\le t}\); composing it with \(\mathsf{H}_{\le t}(\omega)\) produces a terminal view (density operator):
\[
\sigma^{(t,\mathsf{V})}_\omega\ :=\ \mathrm{View}_t(\omega;\mathsf{V})\ \in\ \mathsf{D}(\mathcal{H}_{\mathsf{V}}),
\]
where \(\mathcal{H}_{\mathsf{V}}\) is the verifier's terminal register.
\par

\begin{definition}[Horizon statistical zero-knowledge (HSZK)]\label{def:discussion-hszk}
Fix \(t\ge 0\) and \(\varepsilon\in[0,1]\).
We say the horizon interface \(\mathsf{H}_{\le t}\) satisfies \textbf{horizon statistical zero-knowledge} at error \(\varepsilon\) (denoted \(\mathrm{HSZK}_\varepsilon(t)\)) if for any macroscopic charge \(q\in\mathcal{Q}\), there exists a ``charge-factorized interface'' depending only on \(q\):
\[
\widetilde{\mathsf{H}}_{\le t,q}\ \in\ \mathrm{Comb}\bigl(\mathsf{In}_{\le t}\to \mathsf{Out}_{\le t}\bigr)
\]
such that for any permitted external strategy \(\mathsf{V}_{\le t}\) and any \(\omega\in Q^{-1}(q)\),
\begin{equation}\label{eq:discussion-hszk-view-trace}
\norm{\sigma^{(t,\mathsf{V})}_\omega-\widetilde{\sigma}^{(t,\mathsf{V})}_q}_1\ \le\ \varepsilon,
\qquad
\widetilde{\sigma}^{(t,\mathsf{V})}_q:=\mathrm{View}_t(\widetilde{\mathsf{H}}_{\le t,q};\mathsf{V}).
\end{equation}
Here \(\norm{\cdot}_1\) denotes the trace norm.
\end{definition}

\begin{definition}[Diamond norm distance of combs]\label{def:discussion-comb-diamond}
For any two same-type combs \(X,Y\in \mathrm{Comb}(\mathsf{In}_{\le t}\to\mathsf{Out}_{\le t})\), define their interactive distinguishing norm (diamond norm/strategy norm) as
\begin{equation}\label{eq:discussion-comb-diamond-def}
\boxed{\;
\norm{X-Y}_\diamond
:=\sup_{\mathsf{V}_{\le t}}\ \norm{\mathrm{View}_t(X;\mathsf{V})-\mathrm{View}_t(Y;\mathsf{V})}_1\;}
\end{equation}
where the supremum is taken over all permitted external strategy combs.
\end{definition}

\begin{theorem}[Categorical characterization and diamond norm criterion for HSZK]\label{thm:discussion-hszk-iff-diamond}
Under the setting of Definitions \ref{def:discussion-hszk}--\ref{def:discussion-comb-diamond},
\(\mathrm{HSZK}_\varepsilon(t)\)
holds if and only if for any \(q\in\mathcal{Q}\) there exists \(\widetilde{\mathsf{H}}_{\le t,q}\) such that
\begin{equation}\label{eq:discussion-hszk-diamond}
\boxed{\;
\sup_{\omega\in Q^{-1}(q)}\ \norm{\mathsf{H}_{\le t}(\omega)-\widetilde{\mathsf{H}}_{\le t,q}}_\diamond\ \le\ \varepsilon\;.}
\end{equation}
In this sense, the physical statement ``the horizon does not leak microstate information'' is rigorously promoted to an interactive norm criterion: within a fixed charge sector, the horizon comb family concentrates in the \(\norm{\cdot}_\diamond\) sense to an \(\varepsilon\)-neighborhood of a single charge interface \(\widetilde{\mathsf{H}}_{\le t,q}\).
\end{theorem}

\begin{proof}
By Definition \ref{def:discussion-comb-diamond}, for any fixed \((\omega,q)\) and any candidate \(\widetilde{\mathsf{H}}_{\le t,q}\), the inequality
\(
\norm{\mathsf{H}_{\le t}(\omega)-\widetilde{\mathsf{H}}_{\le t,q}}_\diamond\le \varepsilon
\)
is equivalent to \eqref{eq:discussion-hszk-view-trace} holding for all \(\mathsf{V}_{\le t}\).
Taking the supremum over \(\omega\in Q^{-1}(q)\) yields \eqref{eq:discussion-hszk-diamond}, which is consistent with the quantifier structure of \(\mathrm{HSZK}_\varepsilon(t)\). This completes the proof.
\end{proof}

\begin{remark}[Canonical simulator: microcanonical/maximum entropy representative and charge symmetry center]\label{rem:discussion-canonical-simulator-microcanonical}
The ``charge-factorized interface'' in Definition \ref{def:discussion-hszk} is not required to be unique.
On a given charge sector \(q\), a natural canonical choice is the maximum entropy representative (microcanonical average state):
\[
\bar\rho_q:=\frac{I_q}{\dim\mathcal{H}_q},
\]
where \(\mathcal{H}_q\) is the effective subspace of the charge sector.
If \(\omega\mapsto \mathsf{H}_{\le t}(\omega)\) is affine on states (physical CPTP linearity) and covariant with respect to symmetry within the sector, then
\(\widetilde{\mathsf{H}}_{\le t,q}:=\mathsf{H}_{\le t}(\bar\rho_q)\)
gives a charge-symmetric canonical simulator.
Under this covariant caliber, \(\bar\rho_q\) is the unique symmetric fixed point within the sector, and hence the most natural center choice in the ``charge-symmetric interface class''; the resulting simulator serves as the unified entry point for subsequently introducing spectral statistics into \(\varepsilon\) error control.
\end{remark}

\begin{proposition}[Equivalent characterization of zero-knowledge error and extractable information: Holevo leakage bound and reverse bound]\label{prop:discussion-hszk-holevo}
Fix \(t,q\) and let \(\{\sigma^{(t,\mathsf{V})}_\omega\}_{\omega\in Q^{-1}(q)}\) be the view state family produced by some external strategy \(\mathsf{V}_{\le t}\).
If \(\mathrm{HSZK}_\varepsilon(t)\) holds, then there exists \(\widetilde{\sigma}^{(t,\mathsf{V})}_q\) such that for all \(\omega\in Q^{-1}(q)\), \(\norm{\sigma^{(t,\mathsf{V})}_\omega-\widetilde{\sigma}^{(t,\mathsf{V})}_q}_1\le \varepsilon\).
Then for any distribution \(p\) (supported on \(Q^{-1}(q)\)), the Holevo information satisfies
\begin{equation}\label{eq:discussion-holevo-upper}
\boxed{\;
\chi\bigl(\{p(\omega),\sigma^{(t,\mathsf{V})}_\omega\}\bigr)
\ \le\
2\varepsilon\,\log d_{\mathsf{V}}+2\,h(\varepsilon)\;}
\end{equation}
where \(d_{\mathsf{V}}:=\dim\mathcal{H}_{\mathsf{V}}\) and \(h\) is the binary entropy function.
\par
Conversely, if for all permitted \(\mathsf{V}_{\le t}\) and all distributions \(p\subseteq Q^{-1}(q)\) we have \(\chi\le \delta\),
then for any \(\omega,\omega'\in Q^{-1}(q)\),
\begin{equation}\label{eq:discussion-holevo-reverse}
\boxed{\;
\norm{\sigma^{(t,\mathsf{V})}_\omega-\sigma^{(t,\mathsf{V})}_{\omega'}}_1
\ \le\
\sqrt{8(\ln 2)\,\delta}\;}
\end{equation}
so there exists \(\varepsilon=O(\sqrt{\delta})\) such that \(\mathrm{HSZK}_\varepsilon(t)\) holds (up to constant factors).
Therefore, ``horizon zero-knowledge'' and ``externally extractable microstate information tending to zero'' are quantitatively equivalent characterizations.
\end{proposition}

\begin{proof}
The upper bound \eqref{eq:discussion-holevo-upper} is a continuity estimate for Holevo information: when a family of states lies within the same \(\varepsilon\)-ball in trace distance, their distinguishable information is controlled by \(O(\varepsilon\log d+h(\varepsilon))\) (derivable from the Alicki--Fannes/Audenaert-type inequality).
\par
For the reverse bound \eqref{eq:discussion-holevo-reverse}, take the two-point distribution \(p=\frac12(\delta_\omega+\delta_{\omega'})\).
Then the Holevo information is
\(
\chi=\frac12 D(\sigma_\omega\Vert \bar\sigma)+\frac12 D(\sigma_{\omega'}\Vert \bar\sigma)
\)
where \(\bar\sigma=\frac12(\sigma_\omega+\sigma_{\omega'})\).
By the quantum Pinsker inequality \(D(\rho\Vert\sigma)\ge \frac{1}{2\ln 2}\norm{\rho-\sigma}_1^2\), and noting
\(\sigma_\omega-\bar\sigma=\frac12(\sigma_\omega-\sigma_{\omega'})\), we get
\(
\chi\ge \frac{1}{8\ln 2}\norm{\sigma_\omega-\sigma_{\omega'}}_1^2
\), yielding \eqref{eq:discussion-holevo-reverse}.
Finally, taking a common reference point \(\omega_0\) for all \(\omega\) and setting \(\widetilde{\sigma}:=\sigma_{\omega_0}\), we obtain \(\mathrm{HSZK}_{O(\sqrt\delta)}\) (at the constant factor level). This completes the proof.
\end{proof}

\paragraph{Approximate 2-design \texorpdfstring{$\Rightarrow$}{=>} decoupling \texorpdfstring{$\Rightarrow$}{=>} HSZK}
To strictly translate ``chaotic scrambling'' into a zero-knowledge error upper bound, we employ the decoupling theorem for approximate 2-designs (\cite{SzehrDupuisTomamichelRenner2013DecouplingTwoDesigns}).

\begin{theorem}[Approximate 2-design decoupling implies horizon zero-knowledge (interactive observable upper bound)]\label{thm:discussion-2design-decoupling-hszk}
Let \(A\) denote the system's microscopic degrees of freedom and \(R\) the reference system (which may absorb the verifier's private register to cover the interactive worst case).
Suppose the effective evolution over duration \(t\) consists of two steps: first apply a random unitary \(U\sim\mathcal{E}_t\) to \(A\), then apply a fixed CPTP process \(T:A\to B\) (radiation release/externally accessible reduction), and denote the initial state as \(\rho_{AR}\).
If \(\mathcal{E}_t\) is a \(\delta\)-approximate unitary 2-design, then there exists a marginal state \(\omega_B\) depending only on \(T\) such that
\begin{equation}\label{eq:discussion-decoupling-bound}
\mathbb{E}_{U\sim\mathcal{E}_t}\,
\norm{\,T\bigl((U\otimes I_R)\rho_{AR}(U^\dagger\otimes I_R)\bigr)-\omega_B\otimes\rho_R\,}_1
\ \le\
\mathrm{poly}(\dim A)\cdot \sqrt{\delta}\cdot 2^{-\frac12\Xi_{\min}},
\end{equation}
where \(\Xi_{\min}\) is a negative exponent determined by the min-entropy data of \(\rho_{AR}\) and \(T\) (see \cite{SzehrDupuisTomamichelRenner2013DecouplingTwoDesigns} for the precise expression).
\par
Interpreting \eqref{eq:discussion-decoupling-bound}: at the time scale where the 2-design threshold is reached, the externally visible marginal and the reference system (i.e., microstate label) are approximately tensor-factored in trace distance, so there exists a simulated output \(\omega_B\) depending only on the macroscopic sector, thereby inducing \(\mathrm{HSZK}_{\varepsilon(t)}(t)\) (where \(\varepsilon(t)\) is explicitly controlled by the right side).
\end{theorem}

\paragraph{Spectral form: SFF/frame potential \texorpdfstring{$\Rightarrow$}{=>} approximate design threshold \texorpdfstring{$\Rightarrow$}{=>} HSZK}
The spectral form factor (SFF) and frame potential are computable diagnostic quantities for ``proximity to the Haar second moment'' (see \cite{Liu2018SpectralFormFactors} and related design literature).

\begin{definition}[Spectral form factor (SFF)]\label{def:discussion-sff}
Given a Hamiltonian \(H\) and partition function \(Z(\beta)=\Tr(e^{-\beta H})\), define the standard normalized spectral form factor:
\[
\mathrm{SFF}(\beta,t):=\left|\frac{Z(\beta+\mathrm{i}t)}{Z(\beta)}\right|^2.
\]
\end{definition}

\begin{proposition}[Computable sufficient condition for HSZK from frame potential/spectral diagnostics]\label{prop:discussion-framepotential-sff-to-hszk}
Let \(\mathcal{E}_t\) be the time-smoothed evolution ensemble, \(F^{(2)}(\mathcal{E}_t)\) its second-order frame potential, and \(F^{(2)}_{\mathrm{Haar}}\) the Haar value.
If there exists a function \(\eta(t)\) such that
\[
F^{(2)}(\mathcal{E}_t)-F^{(2)}_{\mathrm{Haar}}\ \le\ \eta(t),
\]
then \(\mathcal{E}_t\) is a \(\delta(t)\)-approximate 2-design (where \(\delta(t)\) is explicitly controlled by \(\eta(t)\)), and by Theorem \ref{thm:discussion-2design-decoupling-hszk} we obtain
\(\mathrm{HSZK}_{\varepsilon(t)}(t)\), with \(\varepsilon(t)\) explicitly controlled by \(\eta(t)\) (and min-entropy data).
\par
In many standard spectral models, \(\eta(t)\) can be controlled by some smoothed upper bound of \(\mathrm{SFF}(\beta,t)\) or its higher moments; thus ``onset/failure of HSZK'' can be reduced to the quantitative determination of when the SFF enters the ramp--plateau region.
\end{proposition}

\paragraph{\texorpdfstring{$\zeta$}{zeta}-spectral model: zeros as the unique spectral input to zero-knowledge error}
Consider the truncated energy spectrum \(E_n=\log n\) (\(n\le N\)).
Let
\(
Z_N(s)=\sum_{n\le N}n^{-s}
\), so that for \(\beta>0\),
\(
Z_N(\beta+\mathrm{i}t)=\sum_{n\le N}e^{-(\beta+\mathrm{i}t)\log n}
\)
is the truncated partition function of this spectral model.
Its SFF thus has the closed form
\[
\mathrm{SFF}_N(\beta,t)=\left|\frac{Z_N(\beta+\mathrm{i}t)}{Z_N(\beta)}\right|^2.
\]
Under appropriate regularization and limiting caliber, the analytic structure of \(Z_N(s)\) coincides with that of \(\zeta(s)\), thereby directly aligning the dip--ramp--plateau behavior of the SFF with the modulus behavior of \(\zeta\) (see \cite{BasuDasKrishnan2025AnalyticZetaRamp}).

\begin{proposition}[Explicit zero-knowledge control law in the \texorpdfstring{$\zeta$}{zeta}-spectral model: zeros enter the error bound as the unique input]\label{prop:discussion-zeta-spectral-hszk}
Under the ``zero-knowledge--SFF control law'' caliber of Proposition \ref{prop:discussion-framepotential-sff-to-hszk},
if \(\eta(t)\) can be controlled by some smoothed upper bound of \(\mathrm{SFF}_N(\beta,t)\), then the error upper bound of \(\mathrm{HSZK}_{\varepsilon(t)}(t)\) can be controlled by a function determined solely by \(|Z_N(\beta+\mathrm{i}t)|\) (and its smoothed average).
In the analytic continuation caliber as \(N\to\infty\), this is equivalent to being controlled solely by \(|\zeta(\beta+\mathrm{i}t)|\).
\par
Since the modulus behavior of \(\zeta\) in analytic number theory is ultimately governed by its zero set (Hadamard decomposition/explicit formulas), in this spectral model:
Riemann zeros appear not as an analogical rhetoric, but as the unique spectral input quantity determining the time behavior of zero-knowledge error.
\end{proposition}

\begin{definition}[Zero-knowledge mixing time]\label{def:discussion-zk-mixing-time}
Define the zero-knowledge mixing time:
\[
t_{\mathrm{ZK}}(\varepsilon):=\inf\{\,t:\ \mathrm{HSZK}_{\varepsilon}(t)\ \text{holds}\,\}.
\]
\end{definition}

\begin{corollary}[Identification of zero-knowledge mixing time with Thouless time (verifiable time scale)]\label{cor:discussion-zk-vs-thouless}
If there exists a time scale \(t_{\mathrm{Th}}\) such that for \(t\ge t_{\mathrm{Th}}\), the spectral diagnostic (SFF or frame potential) enters a stable ramp and reaches the approximate 2-design threshold, then by Proposition \ref{prop:discussion-framepotential-sff-to-hszk} and Theorem \ref{thm:discussion-2design-decoupling-hszk},
\[
t_{\mathrm{ZK}}(\varepsilon)\ \le\ t_{\mathrm{Th}}+O_\varepsilon(1),
\]
where \(O_\varepsilon(1)\) depends on the threshold and smoothing scheme.
For the \(\zeta\)-spectral model, existing analytic work provides evidence that its Thouless time is \(O(1)\) (in its canonical normalization, see \cite{BasuDasKrishnan2025AnalyticZetaRamp}), yielding the co-order conclusion \(t_{\mathrm{ZK}}(\varepsilon)=O(1)\).
\end{corollary}

\begin{conjecture}[Spectral rigidity conjecture for horizon zero-knowledge: HSZK \texorpdfstring{$\Rightarrow$}{=>} sine-kernel correlations]\label{conj:discussion-horizon-zk-rigidity}
Suppose there exists a family of horizon interfaces \(\{\mathsf{H}^{(d)}_{\le t}\}\) with system dimension \(d\to\infty\), and a time scale \(t(d)=\mathrm{poly}(\log d)\) such that
\(\mathrm{HSZK}_{o(1)}(t(d))\)
holds uniformly for all polynomial-complexity external interactive strategies.
Then it is conjectured that the unfolded energy levels of the corresponding effective Hamiltonian at the local unfolding scale must enter the random matrix level repulsion and sine-kernel two-point correlation universality class;
in other words, if the local statistics significantly deviate from the Montgomery-type pair correlation (or its generalizations), then there exists a family of polynomial-complexity interactive verifiers that can construct an observable ``knowledge leakage'' distinguisher, thereby negating \(\mathrm{HSZK}_{o(1)}\).
This conjecture promotes ``zero-knowledge simulability'' to a spectral universality principle, and provides an information-theoretic refutation interface for the pair correlation target in number theory (see \cite{GoldstonLeeSchettlerSuriajaya2025PairCorrelationConjecture} for advances on PCC).
\end{conjecture}

\begin{remark}[Near-horizon--zero realization screening principle: constraining zero realization mechanisms via HSZK]\label{rem:discussion-zk-filter-zero-realization}
Proposals for near-horizon dynamical realizations of Riemann zeros (e.g., through matching near-horizon scalar field quasinormal mode frequencies with nontrivial zeros) can provide realization candidates for the zero set at the spectral level (see models of the type in \cite{Giri2010PhysicalRealizationRiemannZerosBlackHole}).
This subsection provides an additional falsifiable screening principle: if one simultaneously requires that ``external interactions under fixed macroscopic charge do not leak microstates'' (i.e., the strong form of \(\mathrm{HSZK}_\varepsilon\)), then the effective evolution must induce approximate 2-design and decoupling in the interactive sense, so its SFF must enter the ramp--plateau universality class and satisfy quantitative thresholds consistent with \(\varepsilon(t)\).
Therefore, a ``zero realization'' is no longer required merely to match some spectral set, but must simultaneously pass the dual constraints of zero-knowledge simulability and spectral form factor universality; this principle advances formal analogy to an auditable interface condition.
\end{remark}

\endinput
